\begin{description}
\item[(1)] This can happen if the sidereal time is 00:00:00 just after the
  midnight and just before the midnight on the same day. It means that the Sun
  should be near the Autumn (September) Equinox. Approximate R.A. of the Sun is
  $18^{\mathrm{h}}$.
  % sic: the official solution says 18 h; at the September equinox the Sun's
  % R.A. is 12 h.
  \hfill(3 points)
\item[(2)] At $0^{\mathrm{h}}$, 1st January, 2018 (JD\,2458119.5), the mean
  sidereal time $\mathrm{GMST}_0$ was $6.706^{\mathrm{h}}$, i.e.
  $\mathrm{GMST}_0 = 100.59^\circ$. Since one mean solar day is longer than one
  mean sidereal day by $0.9856^\circ$ (about 4 minutes), the GMST will increase
  by $\Delta = 0.9856^\circ$ at next midnight on 2nd January, 2018. Therefore,
  we can first calculate how many days it will take for GMST to reach
  $360^\circ$, i.e.
  \[
    T = \frac{360^\circ - \mathrm{GMST}_0}{\Delta}
      = \frac{360^\circ - 100.59^\circ}{0.9856^\circ} = 263.20
  \]
  \hfill(3 points)

  i.e.\ after completion of 263 days, the GMST at midnight will be
  $359.80^\circ$.

  Thus, GMST will be 00:00:00 on the 264th day at 00:00:48 and then again at
  23:56:52 on the same day. \hfill(2 points)

  The 264th day in the calendar of 2018 is 21st September 2018.
  \hfill(2 points)
\end{description}
